\documentclass[12pt,a4paper]{article}
\usepackage[utf8]{inputenc}
\usepackage{xepersian}
\usepackage{graphicx}
\usepackage{hyperref}
\usepackage{enumitem}

\settextfont{XB Niloofar}
\setdigitfont{XB Niloofar}

\title{پایگاه داده}
\author{}
\date{}

\begin{document}

\maketitle

با رشد تعداد کاربران، یک سرور کافی نیست و به چندین سرور نیاز داریم: یکی برای ترافیک وب/موبایل و دیگری برای پایگاه داده (شکل ۱-۳). جداسازی ترافیک وب/موبایل\footnote{وب تایر - Web Tier} و سرورهای پایگاه داده\footnote{دیتا تایر - Data Tier} امکان مقیاس‌پذیری مستقل آن‌ها را فراهم می‌کند.

\section{کدام پایگاه داده را انتخاب کنیم؟}

می‌توانید بین پایگاه داده‌های رابطه‌ای سنتی و غیررابطه‌ای یکی را انتخاب کنید. بیایید تفاوت‌های آن‌ها را بررسی کنیم.

\subsection{پایگاه داده‌های رابطه‌ای}

به این نوع پایگاه‌داده‌ها سیستم مدیریت پایگاه دادهٔ رابطه‌ای\footnote{RDBMS - Relational Database Management System} یا پایگاه دادهٔ SQL نیز گفته می‌شود. محبوب‌ترین نمونه‌های آن شامل MySQL، Oracle Database، PostgreSQL و غیره هستند.

\begin{itemize}
\item داده‌ها در این پایگاه‌ها به صورت جدول و سطر نمایش و ذخیره می‌شوند.
\item می‌توانید از عملیات جوین\footnote{Join} در SQL برای ارتباط بین جداول مختلف استفاده کنید.
\end{itemize}

\subsection{پایگاه داده‌های غیررابطه‌ای}

این نوع پایگاه‌ها با نام NoSQL نیز شناخته می‌شوند. نمونه‌های معروف آن CouchDB، Neo4j، Cassandra، HBase و Amazon DynamoDB هستند\footnote{منبع ۲}.

این پایگاه‌ها به چهار دسته تقسیم می‌شوند:

\begin{enumerate}
\item ذخیره‌سازی کلید-مقدار\footnote{Key-Value Stores}
\item ذخیره‌سازی گراف\footnote{Graph Stores}  
\item ذخیره‌سازی ستونی\footnote{Column Stores}
\item ذخیره‌سازی سند\footnote{Document Stores}
\end{enumerate}

معمولاً از عملیات جوین در پایگاه‌های غیررابطه‌ای پشتیبانی نمی‌شود.

\section{کدام یک را انتخاب کنیم؟}

برای اکثر توسعه‌دهندگان، پایگاه‌های دادهٔ رابطه‌ای بهترین گزینه هستند، چون بیش از ۴۰ سال سابقه دارند و عملکردشان به‌خوبی اثبات شده است. با این حال، اگر پایگاه‌های رابطه‌ای برای نیازهای خاص شما مناسب نباشند، بررسی گزینه‌های غیررابطه‌ای ضروری است.

\textbf{پایگاه دادهٔ غیررابطه‌ای ممکن است گزینهٔ مناسبی باشد اگر:}

\begin{itemize}
\item برنامهٔ شما نیاز به تأخیر بسیار کم\footnote{Super-Low Latency} دارد.
\item داده‌های شما بدون ساختار هستند یا هیچ رابطه‌ای بین آن‌ها وجود ندارد.
\item فقط نیاز به سریالایز و دیسریالایز داده‌ها (مثل JSON، XML، YAML و غیره) دارید.
\item نیاز به ذخیرهٔ حجم عظیمی از داده‌ها دارید.
\end{itemize}

\end{document}
